\section{An unconditional 269-class certificate for the first quartic reflection field}
\label{sec:cubic-quartic-class-certificate}

This section certifies only the prime-to-$269$ assertion needed for
one quartic class sector.  A provisional class group is used to choose
an integer to test, never as a completeness certificate.  Put
\[
 p=269,\quad t^3=p,\quad z=\zeta_5,\quad
 L=\mathbb Q(t,z),\quad M=L(\sqrt{-3}).
 \tag{QCC.1}
\]
The exact units and finite-field data below, together with outward
rounded analytic intervals, give the following finite certificate.

\begin{theorem}[AC-G129-quartic-class: the first quartic class sector]
\label{AC-G129-quartic-class}
The ideal class group of $L$ has no $269$-torsion.  Consequently the
nontrivial cubic-character part of $\operatorname{Cl}(M)[269]$ is zero,
including its two quartic-character components used for $e=267$.
This does not certify the other degree-twelve fields of an auxiliary
biquadratic host, the $e=133$ field, or H04.4.
\end{theorem}

\begin{proof}
\proofstep{1. Field data without a provisional class-group assumption.}
Let $T=\mathbb Q(t)$ and $C=\mathbb Q(z)$.  Their degrees are three
and four, so they are linearly disjoint.  The polynomial $X^3-p$ is
Eisenstein at $p$.  Its local extension there is tame of degree three
and different exponent two.  At $3$, the congruence
$p\equiv-1\pmod9$ makes $p$ a cube in $\mathbb Q_3$: cubing the
$3$-adic unit group has image the units congruent to $\pm1$ modulo9.
Hence $T\otimes\mathbb Q_3=\mathbb Q_3\times\mathbb Q_3(\sqrt{-3})$
and its discriminant exponent is one.  There is no other possible
ramification, by the polynomial discriminant.  Thus
$D_T=-3p^2$.  The cyclotomic discriminant is $D_C=5^3$.
At every rational prime one of these two integer algebras is etale;
unramified base change commutes with integral closure.  Their tensor
product is therefore the full integer ring of the compositum locally,
and the discriminant tensor formula gives
\[
 (r_1(L),r_2(L))=(0,6),\qquad
 D_L=(3p^2)^4 5^9=4337438179275555809220703125.
 \tag{QCC.2}
\]
The intersection of the $S_3$ splitting field with $C$ is abelian
over $\mathbb Q$, hence is either $\mathbb Q$ or $K$.  The latter
is impossible because $C$ is unramified at3.  Thus the normal closure
has group $S_3\times C_4$, and $L$ is fixed by a reflection in the
first factor.  The normal subgroup generated by
that reflection is $S_3$.  Every abelian subextension of $L/\mathbb Q$
is consequently contained in $C$.  Since a field generated by a
root of unity is abelian, $\mu_L=\mu_C=\mu_{10}$.

\proofstep{2. Exact units and a finite $269$-saturation certificate.}
Use the ordered rational basis $t^a z^b$, $0\le a<3$, $0\le b<4$,
in lexicographic order.  The five units have numerator vectors
\[
\begin{aligned}
 U_1={}&(-1,-1,0,0;0,0,0,0;0,0,0,0),\\
 U_2={}&(270,0,0,0;-739,0,0,0;108,0,0,0),\\
 U_3={}&(-93131838,-243861415,-243861415,-93131838;\\
       &-14428536,-37777370,-37777370,-14428536;\\
       &-2235972,-5852452,-5852452,-2235972),\\
 U_4={}&(123029571,373382222,405083337,174314690;\\
       &-39479967,-52428731,-20952996,11450068;\\
       &3164085,-837481,-6474546,-5956867),\\
 U_5={}&(-405083337,-123029571,-174314690,-373382222;\\
       &20952996,39479967,-11450068,52428731;\\
       &6474546,-3164085,5956867,837481).
\end{aligned}
 \tag{QCC.3}
\]
Divide the last three vectors by3; the first two have denominator1.
This specifies $u_1,\ldots,u_5$ exactly.  The rational multiplication
rule uses only $t^3=269$, $z^5=1$, and $z^4=-1-z-z^2-z^3$.
The certificate forms the twelve multiplication columns of each $u_i$
and computes its characteristic polynomial over $\mathbb Q$.
Each resulting polynomial is monic, has integral coefficients and
constant1; evaluation at $u_i$ gives zero exactly.  Thus both $u_i$
and $u_i^{-1}$ are integral, proving that these are global units.
The complete five polynomials and exact rational verification are in
\texttt{all\_units\_exact\_certificate.json} and
\texttt{verify\_all\_units\_exact.py}; no claim that these are a full
fundamental system is used.

Here is the smaller certificate that matters at269.  The prime
$\ell=40351$ is verified by division through $\lfloor\sqrt\ell\rfloor$.
Its multiplicative group has primitive root3, as checked against the
prime divisors $2,3,5,269$ of $\ell-1$.  Set
$h=3^{150}=10281\pmod\ell$, of order269.  The first five reduction
maps have $(t,z)$ equal respectively to
\[
 \begin{gathered}
 (21085,8631),\ (21085,6215),\ (21085,15186),\\
 (21085,10318),\ (23039,8631).
 \end{gathered}
 \tag{QCC.4}
\]
Cubing and fifth powering check the defining equations, with primitive
fifth roots.  Since $\ell\nmid15p$, these are genuine residue maps of
$\mathcal O_L$, not merely maps of a nonmaximal order.
Write the image of $u_i^{150}$ under map $j$ as $h^{c_{ji}}$.
Direct modular powering gives
\[
 (c_{ji})=
 \begin{pmatrix}
 190&49&187&141&75\\
 79&49&82&75&46\\
 79&49&82&7&141\\
 190&49&187&46&7\\
 190&142&148&73&95
 \end{pmatrix},\qquad
 \det(c_{ji})=-3919610628\equiv208\pmod{269}.
 \tag{QCC.5}
\]
These are homomorphisms $\mathcal O_L^\times\to\mathbb F_{269}$
that annihilate every global $269$th power and every root of unity.
The nonzero determinant proves independence of the five classes.
It also proves multiplicative independence: an integral relation
modulo torsion would have all exponents divisible by269, and division
of that relation by269 can be repeated until a nondivisible exponent
occurs.  This is impossible.  The unit rank is five, so
\[
 U=\mu_L\langle u_1,\ldots,u_5\rangle,\qquad
 m=[\mathcal O_L^\times:U]<\infty,\qquad 269\nmid m.
 \tag{QCC.6}
\]
Indeed, if the finite quotient had an element of order269, its lift
would give a nonzero vector in the kernel of (QCC.5).  Torsion has
order10 and therefore introduces no exception.

\proofstep{3. An explicit lower bound for the true regulator.}
We give a deliberately coarse bound; no unverified table of regulators
is required.  For degree $n$ set
\[
 P_n=\prod_{j=1}^{n}\bigl(4\binom nj+1\bigr),\qquad
 \delta_n=\frac{\log2}{n(P_n+1)}.
 \tag{QCC.7}
\]
A nontorsion unit has a conjugate of modulus greater than
$\exp(\delta_n)$.  To prove this, suppose all its conjugates had
modulus at most that number.  For $1\le k\le P_n+1$, the degree-$n$
characteristic polynomial of its $k$th power has integral $j$th
coefficient of absolute value at most $2\binom nj$.  There are at
most $P_n$ possible polynomials.  Two powers therefore have the same
polynomial.  Comparing their largest root moduli forces that modulus
to be1.  The unit norm then makes all conjugate moduli1.  In this
case the polynomials of all positive powers have bounded integral
coefficients, and there are only finitely many possible roots of
those finitely many polynomials.  Two powers of the original unit
are equal, making it torsion, a contradiction.

Let $r=r_1+r_2-1$ and use the first $r$ weighted archimedean logarithms;
the last coordinate is the negative sum.  If every retained coordinate
has absolute value less than $\delta_n/r$, every unweighted log modulus
has absolute value less than $\delta_n$.  Thus the projected unit
lattice has no nonzero point in that open cube.  The translates of
the half-size cube inject almost everywhere into its real torus:
two points with the same image would have a lattice difference in
the forbidden cube.  Comparing volumes gives
\[
 R_L\ge(\delta_n/r)^r.
 \tag{QCC.8}
\]
For $n=12,r=5$, use $\log2>1/2$ and put
\[
 \begin{gathered}
 P=P_{12}=95338564630654861425388299414125,\\
 R_L>\frac1{[120(P+1)]^5}.
 \end{gathered}
 \tag{QCC.9}
\]
The regulator $R_U$ is the absolute determinant of any five rows of
the six-by-five matrix $2\log|\sigma(u_i)|$.  Representatives for the
six complex places are $t\mapsto269^{1/3}\omega^j$, $z\mapsto z^b$,
with $j=0,1,2$ and $b=1,2$.  Outward rounded evaluation of these exact
expressions proves
\[
 20872891<R_U<20872892,\qquad
 m=R_U/R_L< M_*=20872892[120(P+1)]^5.
 \tag{QCC.10}
\]
All six weighted logarithms are computed; their column sums contain
zero, while the five-by-five determinant is bounded away from zero.
The regulator convention here is the coordinate-projection convention
of Sutherland, Lecture19, printed p.4.

\proofstep{4. The exact analytic expression to be enclosed.}
Let $\rho$ be the standard representation of the $S_3$ splitting
field of $X^3-p$, let $\psi(2)=i$ be primitive modulo5, and let
$\delta=\psi^2$.  The permutation representation on $L$ is
$(1\oplus\rho)\otimes(1\oplus\delta\oplus\psi\oplus\bar\psi)$.
Artin formalism, including inertia invariants at ramified primes,
gives
\[
 \operatorname*{res}_{s=1}\zeta_L(s)
 =L(1,\delta)|L(1,\psi)|^2
   L(1,\rho)L(1,\rho\delta)|L(1,\rho\psi)|^2.
 \tag{QCC.11}
\]
The finite Fourier identity for the quadratic character modulo5 is
$\sum_{a=1}^4\delta(a)e^{2\pi i an/5}=\sqrt5\,\delta(n)$.
Multiply the absolutely convergent identity
$-\log|1-re^{i\theta}|=\sum_{n\ge1}r^n\cos(n\theta)/n$
by $\delta(a)$, sum over $a$, and let $r\uparrow1$.
Dirichlet convergence and Abel's limit theorem justify the limit.
This gives
$L(1,\delta)=-\sum_{a=1}^4\delta(a)\log(2\sin(\pi a/5))/\sqrt5$.
Pairing $a$ with $5-a$ and using
$\sin(2\pi/5)/\sin(\pi/5)=2\cos(\pi/5)=\varphi=(1+\sqrt5)/2$
gives $L(1,\delta)=2\log\varphi/\sqrt5$ directly.
For the odd quartic character, $B_{1,\psi}=(-3-i)/5$.
The primitive Dirichlet functional equation, with Gauss sum of
absolute value $\sqrt5$, gives $|L(1,\psi)|^2=2\pi^2/25$.
Consequently their product in (QCC.11) is exactly
\[
                 \frac{4\pi^2\log\varphi}{25\sqrt5}.
 \tag{QCC.12}
\]
The Dirichlet identities used here are DLMF25.15.3,25.15.5--6,
with $\zeta(0,a)=1/2-a$; they do not use RH.

Write $L(s,\rho)=\sum a_n n^{-s}$.  The coefficients are computed
multiplicatively.  At a prime $\ell\ne3,p$, put
$a_\ell=\#\{x\bmod\ell:x^3=p\}-1$ and
$\eta_\ell=(-3/\ell)$.  Then
\[
 a_{\ell^0}=1,\quad
 a_{\ell^{v+1}}=a_\ell a_{\ell^v}-\eta_\ell a_{\ell^{v-1}},
 \quad a_{3^v}=1,\quad a_{p^v}=0\ (v\ge1).
 \tag{QCC.13}
\]
The roots are either one at inert primes, or three/zero at split
primes, detected by $p^{(\ell-1)/3}\bmod\ell$.
At3 the local cubic Kummer extension is trivial, as proved above;
at$p$ its cubic inertia has no standard invariants.  This proves the
two exceptional Euler factors in (QCC.13).
The twists have coefficients $a_n\delta(n)$ and $a_n\psi(n)$,
including zeros at multiples of5.  Their primitive conductors are
\[
                    Q_0=3p^2=217083,\qquad
                    Q_5=75p^2=5427075.
 \tag{QCC.14}
\]
For this conductor assertion, $\rho=\operatorname{Ind}_{K}^{\mathbb Q}\chi$
for the cubic Kummer character.  Its conductor over $K$ is $(p)$:
it is tame at$p$, trivial at3, and unramified elsewhere.  The norm
pullbacks of $\delta$ and $\psi$ have conductor $(5)$; these nontrivial
characters of $\mathbb F_{25}^\times$ do not cancel the cubic factor
at$p$.  Induction multiplies the conductor norm by $|D_K|=3$.

\proofstep{5. Rapid evaluation with an unspecified but certified root number.}
Here is a direct theta argument fixing every normalization needed
by the computation.  The ring $\mathcal O_K=\mathbb Z[\omega]$ has
class number one.  A character just used is a primitive character
$\chi_f$ modulo $(f)$, with $f=p$ or $5p$, trivial on $\mu_6$.
For the cubic factor the last assertion follows from $p+1$ being
divisible by9; norm characters are trivial on these units automatically.
Extend $\chi_f$ by zero on nonunits and set
\[
 a=\frac{2\pi}{\sqrt3 f}=\frac{2\pi}{\sqrt Q},\qquad
 \theta(t)=\frac16\sum_{\alpha\in\mathcal O_K}
       \chi_f(\alpha)e^{-a|\alpha|^2t}
       =\sum_{n\ge1}b_n e^{-ant}.
 \tag{QCC.15}
\]
The second equality follows by choosing the six generators of every
principal ideal.  The finite Fourier transform of a primitive
multiplicative character is zero at a nonunit frequency, and at a
unit frequency is a constant times its conjugate character.
For the latter assertion substitute the inverse frequency in the
finite sum.  For the former choose a unit on which the character is
nontrivial but which fixes the additive character; primitivity at the
missing prime supplies such a unit.  Orthogonality of additive
characters shows that the absolute value of the constant is $f$.

Poisson summation on $f\mathcal O_K$ now gives
\[
                \theta(t)=\varepsilon t^{-1}\bar\theta(1/t),
                         \qquad |\varepsilon|=1.
 \tag{QCC.16}
\]
Indeed the lattice area is $f^2\sqrt3/2$, its Euclidean dual is
$2/(f\sqrt{-3})\mathcal O_K$, and the Gaussian transform is
$\pi/(at)$ times $\exp(-\pi^2|w|^2/(at))$.  These constants change
$a t$ to $a/t$ and multiply the normalized finite Fourier transform
by $1/(ft)$.  Its Gauss constant has modulus$f$.  More explicitly, write
$d=2/\sqrt{-3}$ and a dual vector as $d\beta/f$.
The residue Fourier sum is $G\,\overline{\chi_f}(\bar\beta)$,
where the fixed factor from $d$ is absorbed in $G$.
Changing the lattice variable $\beta$ to $\bar\beta$ preserves
its norm and gives exactly $\bar\theta(1/t)$.
Thus $\varepsilon=G/f$ proves (QCC.16) without assuming its phase.

Mellin transformation and splitting the integral at1 gives, at$s=1$,
\[
 L(1)=\sum_{n\ge1}\frac{b_n e^{-an}}n
       +\varepsilon a\sum_{n\ge1}\overline{b_n}E_1(an),\qquad
 E_1(x)=\int_1^\infty\frac{e^{-xt}}t\,dt.
 \tag{QCC.17}
\]
All terms and transformed tails decay exponentially; the nonprincipal
finite character has no constant term.  The formula follows first
from the entire Mellin integral and then termwise integration on
$[1,\infty)$.  The interval computation proves $\theta(1)\ne0$,
so (QCC.16) determines the phase rigorously by
\[
                    \varepsilon=\theta(1)/\overline{\theta(1)}.
 \tag{QCC.18}
\]
For real coefficients this gives1.  No guessed sign or rounded
Gauss-sum identity is used for the quartic twist.

\proofstep{6. Absolute tails, interval isolation and the class conclusion.}
Every local coefficient in (QCC.13) is at most $v+1$ in absolute value;
hence $|b_n|\le d(n)\le n$.  With $z_a=e^{-a}$ and
$F=e^{-a(N+1)}$, the omitted absolute tails for the first sum in
(QCC.17), the second sum before multiplication by$a$, and theta are
bounded respectively by
\[
 \frac F{1-z_a},\qquad
 \frac F{a(1-z_a)},\qquad
 \frac{F((N+1)-Nz_a)}{(1-z_a)^2}.
 \tag{QCC.19}
\]
Here $E_1(x)\le e^{-x}/x$.  Each residue subtotal is bounded by
this full positive tail, so splitting by residues modulo5 only
increases the final safe radius.  The certificate uses $N=40000$
for $Q_0$ and $N=200000$ for $Q_5$ at900-bit working precision.
Every elementary, special-function and matrix operation uses Arb
balls; the displayed tail bounds are added with outward rounding.

The analytic class number formula in the convention already fixed is
\[
 h_L R_L=\frac{10\sqrt{D_L}}{(2\pi)^6}
                       \operatorname*{res}_{s=1}\zeta_L(s).
 \tag{QCC.20}
\]
Using (QCC.11)--(QCC.20) and the interval determinant gives
\[
 \left|\frac{h_LR_L}{R_U}-189\right|<8\cdot10^{-203},
 \qquad
 \frac1{M_*}>2\cdot10^{-178}.
 \tag{QCC.21}
\]
The exact rational $M_*$, all ball endpoints, logarithm matrices,
phase intervals and tail radii are preserved in
\texttt{d1\_analytic\_interval.json}; the reproducible verifier is
\texttt{certify\_d1\_analytic.py}.  The final comparison is a strict
ball comparison with $1/M_*$, not an inference from printed digits.

Since $R_U=mR_L$, the enclosed quotient is $h_L/m$.
If $h_L\ne189m$, the nonzero integer $h_L-189m$ would imply
$|h_L/m-189|\ge1/m>1/M_*$, contradicting (QCC.21).
Therefore $h_L=189m$.  Both189 and$m$ are prime to269 by (QCC.6),
so $\operatorname{Cl}(L)[269]=0$.  The equality $m=1$ is neither
needed nor asserted; in particular this proof does not certify the
provisional full class-group structure.

Finally let $s$ be the reflection fixing $L$ in $M$.  Extension and
norm of ideal classes have composites multiplication by2 on
$\operatorname{Cl}(L)$ and $1+s$ on $\operatorname{Cl}(M)$.
Thus the $s$-invariant part of $\operatorname{Cl}(M)[269]$ is zero,
and $s$ acts as$-1$.  All conjugate reflections do likewise.
Their products, the cubic rotations, act as1.  The nontrivial cubic
isotypic part is therefore zero, proving the asserted class-sector
conclusion, without asserting that every component of
$\operatorname{Cl}(M)[269]$ vanishes.
\end{proof}

\paragraph{Source and certification scope.}
The analytic class formula is the theorem on printed p.1 of
A.~Sutherland, MIT18.785 Lecture19 (13 November2019), with weighted
logarithms and the projected regulator defined on printed p.4:
\url{https://math.mit.edu/classes/18.785/2019fa/LectureNotes19.pdf}.
Only those statement and normalization pages are claimed as fresh
source inspection, not a fresh whole proof of that classical theorem.
The Dirichlet functional equation is
\url{https://dlmf.nist.gov/25.15}, equations3,5,6.  The specialized
Hecke theta normalization and rational isolation argument have been
proved here.  The interval runtime is python-flint0.9.0, using its
\texttt{arb.expint} order1, exact-integer inputs and the documented
inclusion arithmetic, at
\url{https://python-flint.readthedocs.io/en/latest/arb.html}.
All finite certificates are separate from the failed or still-running
PARI completeness jobs.  No Lean build, GRH for the nonabelian field,
or assumption about the distribution of cubic zeros enters the proof.

\paragraph{Geometric dictionary and remaining gate.}
The class-sector conclusion concerns unramified abelian covers of
$\operatorname{Spec}\mathcal O_M$ and their cubic-character part,
through the retained class-field and finite-etale-cover dictionary
(SGA1, ExposeV, sections4 and7).  Its reflection argument is exact
prime-to269 restriction and norm, not a purity or positivity theorem
for a nonabelian motive.  The degree-two Artin representation and its
finite characters are weight-zero finite-monodromy sheaves away from
the displayed ramification.  This finite arithmetic certificate
removes one class input in the quartic Kummer calculation; it supplies
no H04.4 window inequality, native deformation host, or new comparison
with an automorphic support by itself.
